Le professeur insère dans le montage de la
Fig.~\ref{fig:svt_national_2016_normale_phys_ex2_3}, en série avec le
conducteur ohmique, une bobine d'inductance $L=\qty{0.1}{\H}$ et de résistance
négligeable.

Après avoir chargé de nouveau et totalement le condensateur, le professeur
bascule l'interrupteur en position~(2), à l'instant $t_{0}=0$.
La Fig.~\ref{fig:svt_national_2016_normale_phys_ex2_5} représente les variations
de la tension $u_{c}(t)$ aux bornes du condensateur et de la tension $u_{R}(t)$
aux bornes du conducteur ohmique.

\begin{figure}[H]
    \centering
    \includegraphics[width=.7\textwidth]{2025_03_17_d0909ada06865003c630g-5(2)}
    \caption{}
    \label{fig:svt_national_2016_normale_phys_ex2_5}
\end{figure}
\begin{enumerate}
    \item Montrer que l'expression de l'énergie totale du circuit à un instant
          $t$ s'écrit~:
          \[
              \mathcal{E}=\frac{1}{2}C\times u_{C}^{2}+
              \frac{1}{2}\frac{L}{R^{2}}\times u_{R}^{2}
          \]
    \item Déterminer la valeur de
          $\Delta\mathcal{E}=\mathcal{E}_{1}-\mathcal{E}_{0}$, la variation de
          l'énergie totale du circuit entre les instants $t_{0}=0$ et
          $t_{1}=\qty{3.5}{\ms}$. Interpréter ce résultat.
\end{enumerate}

